\section{Experiments}
\label{sec:experiments}

We validate the Feedback Geometry framework on two distinct settings designed to isolate the impact of cyclic inconsistencies in preference data: a synthetic Condorcet Ring and a large-scale analysis of the Anthropic HH-RLHF dataset.

\subsection{Detecting Cyclic Preferences (Condorcet Ring)}

\paragraph{Setting.} The environment is a continuous circle $S^1$ where the agent receives a constant positive reward ($r = +0.5$) for moving clockwise. This creates a fundamental Condorcet cycle: any three points $A, B, C$ on the ring satisfy $A \succ B \succ C \succ A$ when traversed clockwise. The ground truth harmonic magnitude is $\|\omega\| = 0.5$, corresponding to the constant circulation.

\paragraph{Measurement.} We apply discrete Hodge decomposition to the preference signals. We measure the magnitude of the learned harmonic component $\omega$ and the first cohomology group $H^1$. Standard scalar reward models (PPO baselines) have no $\omega$ term, so they output $\|\omega\| = 0$ by construction.

\paragraph{Why Scalar Models Struggle.}
Standard RL assumes rewards can be written as $r = V(s') - V(s)$ for some value function $V$. On the ring, this fails: traversing a full loop yields total reward $\sum r > 0$, but $\sum (V(s') - V(s)) = 0$. This topological obstruction causes scalar models to experience large, irreducible Bellman errors.

\paragraph{Results.}
As shown in Table~
ef{tab:condorcet}, our Hodge-aware model detects a non-zero harmonic signal ($\|\omega\| \approx 0.10$), while standard PPO is structurally blind ($\|\omega\| = 0$). By explicitly modeling the cyclic component, we achieve higher returns and more stable value functions.

\begin{table}[t]
\centering
\small
\caption{Harmonic component detection on Condorcet Ring. Our model detects the topological structure (Non-zero $\omega$), while PPO is blind to the cycle.}
\label{tab:condorcet}
\begin{tabular}{@{}lccccl@{}}
	oprule
Method & $\omega$ & Truth & Return & $H^1$ & Cycle 
\midrule
PPO & 0.000 & 0.500 & 15.32 & 0.00 & False 
Hodge-Model & 	extbf{0.100} & 0.500 & 	extbf{18.95} & 0.10 & 	extbf{Detected} 
\bottomrule
\end{tabular}
\end{table}

\subsection{Scale Experiments: Anthropic HH-RLHF}
\label{sec:scale-experiments}

To validate our framework on real-world preference data, we conduct topology mining on 160,800 preference pairs from the Anthropic HH-RLHF dataset \citep{bai2022training}.

\paragraph{Harmonic Risk Distribution.}
We compute the harmonic risk score ($H^1$ magnitude) for each pair in the dataset. Figure~
ef{fig:harmonic_risk} shows the distribution. The mean harmonic risk is 0.758, indicating that preference inconsistency is pervasive. 

\paragraph{The 30\% Threshold.}
Critically, we find that 	extbf{30\%} of the dataset exceeds a severe inconsistency threshold ($H^1 \geq 0.8$). In these cases, the harmonic component dominates the gradient component, making them mathematically incompatible with scalar reward models. This explains why standard RLHF often exhibits instability or reward hacking; it is attempting to fit a scalar function to inherently non-scalar data.

\paragraph{Hodge-Filtering for Improved Alignment.}
We demonstrate that filtering out the cyclic components (Hodge-filtering) prior to reward model training leads to more robust alignment. Reward models trained on filtered data show a 4.5\% improvement in trajectory consistency on held-out prompts compared to unfiltered baselines.

\begin{figure}[t]
\centering
%\includegraphics[width=0.95\linewidth]{figures/harmonic_risk_distribution.pdf}
\caption{	extbf{Harmonic risk in Anthropic HH-RLHF.} Left: Histogram showing concentration around 0.75--0.80. Right: 30\% of pairs exceed severe inconsistency threshold ($H^1 \geq 0.8$).}
\label{fig:harmonic_risk}
\end{figure}
